% ---------------------------------------------------
% Trabajo Final de Grado
% Author: Fabián González Lence <alu0101549491@ull.edu.es>
% Chapter: Resultados
% ----------------------------------------------------

\chapter{Resultados} \label{chap:Results}

Los cinco proyectos desarrollados en el experimento están disponibles de forma pública y accesible: los tres primeros exclusivamente como \textit{frontend} estático, y los dos últimos con \textit{backend} desplegado en la nube. Este capítulo documenta la infraestructura de despliegue adoptada (Sección~\ref{sec:results-deployment}), las métricas de calidad de código obtenidas mediante SonarQube\footnote{\href{https://sonarcloud.io}{SonarQube Cloud}} y el análisis del agente de \ac{IA} (Sección~\ref{sec:results-code-quality}), y una valoración global de los proyectos resultantes (Sección~\ref{sec:results-analysis}).

\section{Despliegue de los proyectos} \label{sec:results-deployment}

Los cinco proyectos se despliegan desde el mismo repositorio monorepo bajo dos estrategias según su naturaleza. Los cinco \textit{frontends} se publican como sitios estáticos en GitHub Pages mediante un flujo de GitHub Actions que construye cada proyecto en secuencia, inyecta la variable \texttt{BASE\_URL} en el fichero \texttt{vite.config.ts} de cada uno para que se sirva bajo su propio subdirectorio, y se genera una página de entrada que enlaza a los cinco subproyectos\footnote{\href{https://alu0101549491.github.io/TFG-Fabian-Gonzalez-Lence/}{Página de entrada a los proyectos}}.\\

Los backends de \ac{CARTO} y \ac{TENNIS} se desplegaron en Render\footnote{\href{https://render.com}{Render}}, con base de datos PostgreSQL alojada en Supabase\footnote{\href{https://supabase.com}{Supabase}}. Para \ac{CARTO} se exploró inicialmente Railway\footnote{\href{https://railway.app}{Railway}}, pero las restricciones de su plan gratuito y la incompatibilidad con la estructura monorepo motivaron su cambio. Ambos servicios se definen en el fichero \texttt{render.yaml} de la raíz del repositorio, centralizando la gestión de infraestructura en el mismo. Durante la migración de \ac{CARTO} a Supabase, las credenciales de conexión de la base de datos fueron expuestas accidentalmente en el repositorio público a causa de un error cometido por el agente de GitHub Copilot, estas fueron detectadas por GitGuardian\footnote{\href{https://www.gitguardian.com}{GitGuardian}} y rotadas de inmediato.

\section{Métricas de calidad del código} \label{sec:results-code-quality}

La calidad del código se evaluó mediante dos herramientas complementarias. La primera es SonarQube Cloud\footnote{\href{https://sonarcloud.io/organizations/alu0101549491}{Organización del \texttt{TFG} en SonarQube Cloud}}, integrado en el repositorio vía GitHub Actions con un fichero \texttt{sonar-project.properties} por proyecto. La segunda es un agente personalizado de GitHub Copilot denominado \textit{Code Quality Agent}, generado mediante \textit{meta-prompting} con Claude Opus 4.7, debido a su gran ventana de contexto. El agente aplica el estándar ISO/IEC 25010 \cite{ISO25010:2023} recorriendo las mismas características de calidad que SonarCloud (fiabilidad, seguridad, mantenibilidad, etc.), con el objetivo de contrastar ambas evaluaciones. Las Tablas~\ref{tab:sonarqube-issues} y~\ref{tab:sonarqube-ratings} recogen las métricas obtenidas directamente de la \ac{API} de SonarCloud. La cobertura de \ac{CARTO} y \ac{TENNIS} figura como no disponible porque ambos proyectos utilizan exclusivamente pruebas \textit{end-to-end} con Playwright, cuya cobertura no es compatible con el formato LCOV esperado por SonarCloud.

\begin{table}[H]
\centering
\small
\renewcommand{\arraystretch}{1.3}
\begin{tabular}{|>{\raggedright\arraybackslash}p{2.6cm}|>{\centering\arraybackslash}p{1.4cm}|>{\centering\arraybackslash}p{1.6cm}|>{\centering\arraybackslash}p{2cm}|>{\centering\arraybackslash}p{2cm}|}
\hline
\textbf{Proyecto} & \textbf{Bugs} & \textbf{Vulner.} & \textbf{\textit{Code smells}} & \textbf{Cobertura} \\
\hline
\ac{HANGMAN} & 0  & 0 & 20  & 95,8\,\% \\
\hline
\ac{PLAYER}  & 0  & 0 & 81  & 83,3\,\% \\
\hline
\ac{BALATRO} & 4  & 2 & 91  & 72,0\,\% \\
\hline
\ac{CARTO}   & 7  & 8 & 440 & ---      \\
\hline
\ac{TENNIS}  & 61 & 0 & 561 & ---      \\
\hline
\end{tabular}
\caption{Incidencias y cobertura de tests según SonarCloud por proyecto.}
\label{tab:sonarqube-issues}
\end{table}

\begin{table}[H]
\centering
\small
\renewcommand{\arraystretch}{1.3}
\begin{tabular}{|>{\raggedright\arraybackslash}p{2.6cm}|>{\centering\arraybackslash}p{2cm}|>{\centering\arraybackslash}p{2cm}|>{\centering\arraybackslash}p{2cm}|>{\centering\arraybackslash}p{2.2cm}|}
\hline
\textbf{Proyecto} & \textbf{Fiabilidad} & \textbf{Seguridad} & \textbf{Mant.} & \textbf{\textit{Quality Gate}} \\
\hline
\ac{HANGMAN} & A & A & A & \texttt{OK} \\
\hline
\ac{PLAYER}  & A & A & A & \texttt{OK} \\
\hline
\ac{BALATRO} & B & B & A & \texttt{OK} \\
\hline
\ac{CARTO}   & C & C & A & \texttt{OK} \\
\hline
\ac{TENNIS}  & E & A & A & \texttt{ERROR} \\
\hline
\end{tabular}
\caption{Calificaciones de calidad según SonarCloud. Fiabilidad, Seguridad y Mantenibilidad se expresan en escala A--E.}
\label{tab:sonarqube-ratings}
\end{table}

La única excepción al \textit{Quality Gate} es \ac{TENNIS}: sus 61 bugs concentrados en el módulo de gestión de partidos resultan en fiabilidad E. Los bugs de seguridad de \ac{CARTO} corresponden al patrón de almacenamiento inseguro de JWT \cite{ISO25010:2023}, aunque el proyecto supera el umbral configurado. La Tabla~\ref{tab:quality-results} recoge el análisis complementario realizado por el agente de \ac{IA} según ISO/IEC 25010.

\begin{table}[H]
\centering
\small
\renewcommand{\arraystretch}{1.3}
\begin{tabular}{|>{\raggedright\arraybackslash}p{1.8cm}|>{\centering\arraybackslash}p{1.5cm}|>{\centering\arraybackslash}p{1.5cm}|>{\centering\arraybackslash}p{1.5cm}|>{\centering\arraybackslash}p{1.5cm}|>{\centering\arraybackslash}p{1.5cm}|>{\centering\arraybackslash}p{2cm}|}
\hline
\textbf{Proyecto} & \textbf{Bloq.} & \textbf{Crít.} & \textbf{Mayor} & \textbf{Menor} & \textbf{Total} & \textbf{Valoración} \\
\hline
\ac{HANGMAN} & 0 & 0 & 1  & 5  & 6   & Bueno \\
\hline
\ac{PLAYER}  & 1 & 2 & 9  & 12 & 24  & Mejorable \\
\hline
\ac{BALATRO} & 0 & 3 & 9  & 16 & 28  & Mejorable \\
\hline
\ac{CARTO}   & 5 & 6 & 13 & 6  & 30  & Crítico \\
\hline
\ac{TENNIS}  & 0 & 5 & 12 & 8  & 25  & Crítico \\
\hline
\textbf{Total} & \textbf{6} & \textbf{16} & \textbf{44} & \textbf{47} & \textbf{113} & --- \\
\hline
\end{tabular}
\caption{Recuento de incidencias y valoración global por proyecto según el análisis ISO/IEC 25010 del agente de \ac{IA}.}
\label{tab:quality-results}
\end{table}

Ambas herramientas coinciden en la tendencia general: la calidad decrece a medida que crece la complejidad del proyecto, con \ac{HANGMAN} como referencia positiva y \ac{CARTO} y \ac{TENNIS} como los que acumulan más deuda. Los patrones más destacables son la ausencia de mecanismos de captura de errores en el nivel raíz de los componentes en todos los \textit{frontends}, el uso recurrente de \texttt{as any} a partir de \ac{PLAYER}, y la cobertura de pruebas insuficiente en los proyectos de mayor escala.

\section{Análisis de los proyectos resultantes} \label{sec:results-analysis}

Los cinco proyectos cumplen funcionalmente con las especificaciones de requisitos definidas en el Capítulo~\ref{chap:Design} y están accesibles en producción \cite{URL::TFG-Web}. La degradación de calidad documentada en la Sección~\ref{sec:results-code-quality} es coherente con la literatura sobre generación de código con \ac{LLMs} revisada en el Capítulo~\ref{chap:RelatedWork}: la calidad decrece a medida que crece el tamaño del contexto necesario, y la progresión de escala entre proyectos reproduce ese patrón con precisión. Las incidencias de mantenibilidad —complejidad ciclomática, funciones extensas, uso recurrente de código cuestionable como el \texttt{as any}— son sistemáticas y predecibles, mientras que las de seguridad y fiabilidad —almacenamiento inseguro de credenciales, ausencia de límites de error en los árboles de componentes— solo emergen de una comprensión global del diseño del sistema, algo que los modelos actuales gestionan de forma inconsistente. Las implicaciones de este comportamiento se desarrollan en el Capítulo~\ref{chap:Conclusiones}.